\chapter{免疫系统：识别、反应与记忆}

\begin{kaichang}
翻开你的疫苗接种本——就是记录着你打过乙肝、百白破、麻腮风疫苗的那个小本子。再看一眼你的左手臂，很多人肩膀上有一个小小的圆形疤痕，那是出生不久打卡介苗留下的。

现在做个小游戏。一针疫苗打进去，药液只有零点几毫升，里面的东西几天之内就被身体清理得干干净净。可是它的“效果”却能持续几年甚至几十年。药液早就没了，保护却还在——\textbf{保护你的到底是什么}？它藏在身体的哪个角落？

再猜第二个问题：被纸划破的手指，第一天有点红，第二天肿起来、摸上去热热的，过几天结痂脱落，完好如初。没有人指挥，没有图纸，这场修复是谁组织的？先写下你的猜想，这一章结束时我们回来对答案。
\end{kaichang}

\section{看得见的战场：红、肿、热、痛}

先只看能观察到的东西。皮肤破了或者被蚊子叮了一口之后，那一小块地方会出现四个变化：发红、肿胀、发热、疼痛。这套反应有个专门的名字——炎症。它不挑敌人：细菌、病毒、木刺、蚊子的口水，甚至一颗嵌进皮肤的刺，都能触发几乎一样的四个症状。

再看另一类现象：生病。你得了普通感冒，流鼻涕、嗓子痛、可能发烧（第 60 章讲过，发烧是体温调定点被主动上调，不是身体失控），难受三五天，然后自己好了——大多数感冒并没有“治感冒的药”，市面上的感冒药只是让你舒服一点，真正收拾残局的是身体自己。

还有一类更奇怪的现象：\textbf{第二次}。得过一次水痘的人，几乎一辈子不会再得水痘；照顾水痘病人的爸妈安然无恙，因为他们小时候得过了。身体好像“记得”这个敌人。可它记得的方式很挑剔——得过水痘，照样会感冒、会得腮腺炎。记忆是精确的，一种敌人一份档案。

这三个现象——千篇一律的炎症、自己痊愈的感冒、精确到“一种一份”的记忆——就是这一章要解释的全部。下面把镜头一层层推进去。

\begin{shiyan}
\textbf{观察一次炎症的全过程}（只做观察，不制造伤口）。材料：一支笔、一个日历或笔记本。步骤：等下一次你被蚊子叮了、或者不小心有了浅浅的擦伤，从当天开始，每天固定时间观察一次，记录四个指标：红不红、肿不肿（和旁边正常皮肤比）、摸上去热不热、按一下痛不痛，画一张四行的小表格，连记五到七天。观察点：四个症状哪天最重？是同时消退还是有先有后？结痂是什么颜色、第几天脱落？安全提示：\textbf{绝对不要故意弄伤自己来“做实验”}；不要抠结痂、不要挤肿包；如果红肿范围不断扩大、疼痛加重或发烧，说明可能感染了，马上告诉家长并就医。
\end{shiyan}

\section{第一道防线：城墙、护城河与巡逻队}

把镜头推进到皮肤的表面。在第 59 章，我们把身体看成细胞的合作社会；这个社会的第一道边防不是细胞，而是\textbf{结构本身}。

皮肤最外层是几层已经角化的死细胞，像铺得密不透风的瓦片，绝大多数微生物根本钻不过去。皮肤表面还铺着一层汗液和皮脂混合的“酸膜”，pH 大约 $5.5$，很多细菌在这个酸度下活不好。眼泪、唾液里有一种叫溶菌酶的蛋白质——它本质上是一种酶（第 49 章），能催化细菌细胞壁中糖链的水解，把城墙里混进来的散兵游勇就地拆碎。胃里的盐酸（pH 约 $1.5$～$2$）则把关消化道入口，随食物进来的微生物大多死在这道“化学护城河”里。

万一敌人真的突破了城墙呢？第二梯队上岗：\textbf{先天免疫}。血液和组织里有一批巡逻的“大胃王”细胞——巨噬细胞和中性粒细胞。它们的工作方式朴素得惊人：吞。发现可疑颗粒，就伸出细胞膜把它包进去（第 51 章讲过膜的流动与变形），形成一个囊泡，囊泡与装满酶的溶酶体合并，把吞进去的东西酶解成小碎片。

这里有个关键问题：巡逻队怎么知道什么该吞、什么不该吞？它们没有眼睛，靠的还是分子识别。细菌的细胞壁成分、细菌的 DNA、病毒的遗传物质，有一些化学图案是人体细胞\textbf{从来不生产}的。巡逻细胞表面装着一批受体蛋白（第 48 章讲过，蛋白质的形状决定它抓住谁），专门匹配这些“非我”图案。图案一对上，吞噬、警报同时启动。

注意这套系统的特点：它认的是\textbf{类别}，不是个体。它不管来的是哪一株细菌，只认“细菌味儿”的通用图案。所以先天免疫的反应永远千篇一律——这就是炎症红、肿、热、痛总是长一个样的原因：红肿热是局部血管扩张、血流增加、巡逻队增援的结果，痛来自炎症物质对神经的刺激。它是固定程序，快，但不精确。

\section{精准的钥匙：抗原、抗体和淋巴细胞}

固定程序能解释炎症，却解释不了“得过水痘就不再得”的精确记忆。精确，来自另一套系统：\textbf{适应性免疫}。它的主角是两类白细胞——B 细胞和 T 细胞，统称淋巴细胞。

先说敌人这边的概念：\textbf{抗原}。抗原不是一个专门的“坏分子类别”，而是任何能被免疫系统识别为“非我”的分子——通常是病原体表面的某一段蛋白质或多糖。一种病毒表面可能有几十种不同的抗原，每种抗原上又有好几个可供识别的“局部区域”，就像一栋楼有好几扇门。

再说防御这边的武器：\textbf{抗体}。抗体是 B 细胞制造的 Y 形蛋白质，两只“手”（可变区）的形状恰好能与某一种抗原的某个局部严丝合缝地结合——特异性的原理，和第 49 章里酶与底物的形状匹配完全一样：靠氢键、离子键和疏水作用，把两个表面按互补的形状贴在一起（第 22 章的分子间作用力）。抗体本身不“杀死”任何东西——记住这一点——它做的是\textbf{标记}：粘在病毒表面让它没法钻入细胞，或者裹在细菌外面给吞噬细胞递上一面“吃我”的旗子。

两种淋巴细胞分工不同：

\begin{itemize}[itemsep=2pt]
\item \textbf{B 细胞}：每个 B 细胞表面架着几十万份同一种形状的受体（就是还没分泌出去的抗体）。受体匹配到抗原、并得到辅助信号确认后，这个 B 细胞被激活，疯狂增殖、分化成“浆细胞”——一座抗体工厂，每秒钟能分泌上千个抗体分子。
\item \textbf{T 细胞}：分两类。辅助 T 细胞是“指挥部”，负责核实警报、决定开战还是停火，没有它的确认，B 细胞大多不肯全力开火；杀伤 T 细胞则是“内部检查员”，它不看游离的病原体，只看\textbf{你自己的细胞}——每个细胞都会把内部蛋白质的碎片摆到表面“公示”，一旦公示的碎片显示“此细胞已被病毒占领”，杀伤 T 细胞就命令这个细胞自我了断，连同里面的病毒一起清除。
\end{itemize}

这套“公示”制度值得停下来想一想：身体几乎每一个细胞都在不停地把自己的内部样品摆出来接受检查。识别“非我”，靠的不只是认敌人，还包括认\textbf{自己的每一张身份证}。这为后面要讲的“出错”埋下了伏笔。

\begin{qiaoliang}
抗体是蛋白质，蛋白质在 $37\,^{\circ}\mathrm{C}$ 的水里不停地热运动、碰撞。一个抗体要发挥作用，得先“撞”上它要抓的抗原。你的血液里约 $5$ 升，血浆中抗体（以 IgG 为主）的浓度约 $\ud{10}{g/L}$，抗体分子的摩尔质量约 $\ud{1.5\times 10^{5}}{g/mol}$。算一算：
\[\frac{\ud{10}{g/L}\times \ud{5}{L}}{\ud{1.5\times 10^{5}}{g/mol}}\approx \ud{3\times 10^{-4}}{mol},\qquad 3\times 10^{-4}\times 6.02\times 10^{23}\approx 2\times 10^{20}\ \text{个分子}\]
也就是说，你身体里此刻漂浮着大约 $2\times 10^{20}$ 个抗体分子——两千亿亿个巡逻中的“分子标签”，比全世界人口多几万亿倍。一场战役能打赢，靠的不是某一个分子有多英勇，而是\textbf{足够大的数量}遇上\textbf{足够高的匹配精度}。这正是化学的看家本领：用统计和概率解释宏观结果。
\end{qiaoliang}

\section{先有多样，再有选择}

现在轮到最深的那层规则。请你直面一个几乎不可能的问题：你的基因组只有两万来个基因，可世界上的病原体千奇百怪，病毒还在不断变异出新面孔——身体怎么能为\textbf{从未见过的敌人}预先准备好恰好匹配的抗体？“先造好锁等着钥匙出现”，钥匙的种类够用吗？

答案叫\textbf{克隆选择}，它是整个免疫学的核心，逻辑分三步：

\begin{enumerate}[itemsep=2pt]
\item \textbf{先随机生成}。每个 B 细胞在骨髓里成熟时，会把编码抗体可变区的几段基因片段\textbf{随机洗牌重组}（很像洗牌抽牌，细节见本章挑战层）。结果：你体内同时存在上亿个 B 细胞，每个带着一种几乎独一无二的受体，总共能覆盖大约 $10^{11}$ 种不同的形状。注意因果顺序——多样性是在\textbf{没见过任何敌人之前}就生成好的。
\item \textbf{敌人来做选择}。某种抗原进入身体，在淋巴器官里“巡展”。上亿种受体里，碰巧有那么几个 B 细胞的受体能抓住它。抗原的作用就是\textbf{把这几个细胞挑出来}：匹配者被激活。
\item \textbf{被选中的克隆扩张}。激活的 B 细胞开始分裂，一个变两个、两个变四个……几天之内扩增成成千上万一模一样的“克隆”，有的变成分泌抗体的浆细胞投入战斗，有的变成\textbf{记忆细胞}长期驻留。战斗结束后，大部分作战细胞按计划凋亡，记忆细胞留下来——可以活几年甚至几十年。
\end{enumerate}

\begin{center}
\begin{tikzpicture}[scale=0.9]
\foreach \i in {0,1,2,3,4,5} \draw[fill=white] (\i*1.1,2.2) circle (0.28);
\node at (0.6,1.5) {\footnotesize 百万种 B 细胞，各带一种受体};
\draw[->,thick] (3.2,2.2) -- (3.2,1.2);
\node at (4.9,1.7) {\footnotesize 抗原进入，只匹配其中一种};
\draw[fill=gray!30] (3.2,0.7) circle (0.28);
\draw[->,thick] (3.2,0.3) -- (3.2,-0.5);
\draw[fill=gray!30] (2.2,-1.0) circle (0.24);
\draw[fill=gray!30] (3.2,-1.0) circle (0.24);
\draw[fill=gray!30] (4.2,-1.0) circle (0.24);
\node at (7.6,-1.0) {\footnotesize 克隆扩增：浆细胞 $+$ 记忆细胞};
\end{tikzpicture}
\end{center}
\vspace{-1em}
（这是人为的示意图：圆圈不代表细胞的真实外观，只表示“许多种里选中一种”。）

请你再读一遍这三步，有没有觉得眼熟？\textbf{先有多样性，再被环境筛选，选中者扩增}——这完全就是自然选择的逻辑（第 76 章会正式展开），只不过这一次发生在你的身体内部，时间尺度从几十万年压缩到几天。免疫系统和进化，用的是同一套算法。这也解释了为什么这个过程不需要“预知未来”：身体从不猜测敌人长什么样，它只是先把彩票印够，再等中奖号码。

第二次遇到同一种抗原时，记忆细胞直接跳过“大海捞针”的阶段，几小时内就开始扩增，抗体浓度升得又快又猛——往往在症状出现之前就把敌人压了下去。这就是“得过水痘就不再得”的机制：不是水痘病毒不敢再来，而是它每次来都被记忆当场认出来，秒杀在萌芽状态。

\begin{zhengju}
现在这套“先有受体、抗原只做选择”的理论，在 1950 年代却是离经叛道。当时的主流是“指令学说”：抗体像一团橡皮泥，抗原进来了当模板，抗体照着抗原的形状现场折叠出互补的形状。这个理论听起来很合理——省材料，需要才造，多么经济！

1957 年，澳大利亚免疫学家伯内特（F.~Burnet）和丹麦的耶讷（N.~Jerne）等人提出了相反的主张：抗体的形状在抗原出现\textbf{之前}就由每个淋巴细胞随机定好了，抗原不做模板，只做“裁判”——谁碰巧匹配，谁就增殖。这就是克隆选择学说。伯内特为此做了一个可检验的预言：如果学说成立，那么每个淋巴细胞应该只产生一种特异性的抗体；而按指令学说，一个细胞拿到什么模板就该能造什么抗体，没理由只能造一种。

判决性实验在几年后到来：研究者把单个淋巴细胞分离出来单独培养，观察它和它的后代——结果，一个细胞分裂出的整个克隆，确实只认一种抗原，给别的抗原毫无反应。再后来，分子生物学直接找到了“随机洗牌”的机器：编码抗体的基因在 B 细胞发育时被重新拼接，洗牌方式是写在 DNA 层面的（第 67 章以后你会知道 DNA 是怎样存储信息的）。指令学说被放弃，不是因为它“不好听”，而是因为它预言的现象没出现，而对手的预言出现了。伯内特后来因此获得诺贝尔奖，但这个故事的重点不是谁聪明，而是：\textbf{两个都能解释“抗体能匹配抗原”的理论，最终靠一个不一样的可检验预言分了胜负}。
\end{zhengju}

\section{疫苗：一次安全的军事演习}

理解了克隆选择和记忆细胞，疫苗的原理就一句话：\textbf{把敌人的“照片”送给免疫系统，让身体在真敌人到来之前先建好档案}。

疫苗里装的，是病原体的抗原——可能是被杀死的整个病原体，可能是减毒到没法致病的活病原体，也可能只是病原体表面的一小段蛋白质，甚至是一段让身体细胞临时生产这段蛋白质的指令（mRNA 疫苗，原理涉及第 68 章的转录与翻译，这里不展开）。无论哪种，共同点都是：\textbf{提供识别目标，不提供完整的敌人}。

这里必须把一条最常见的误会拆开，它是本章的红线：

\textbf{疫苗不是药物，它替你杀死的病原体是零。}药片里的抗生素直接作用于细菌（第 62 章会讲它带来的耐药性问题）；而疫苗进入身体后，直接杀灭的病原体数量是\textbf{零}——它做的全部事情，是触发一次克隆选择：匹配的 B 细胞被选中、扩增、留下记忆。几个月后真正的病毒入侵时，冲上去作战的是\textbf{你自己的}抗体、你自己的记忆细胞。药液早就分解消失了，留下来的是被训练过的免疫系统。打个比方：疫苗是消防演习，不是替你请来的消防队——演习结束，消防员还是你自己的。

所以，疫苗的效果依赖被接种者的免疫系统正常工作，存在个体差异；它的保护率和持久性要靠\textbf{大规模人群试验与长期监测数据}来评估，不同疫苗差别很大。也正因如此，本章和全书一样，只讲证据和机制，不提供任何个人的接种或诊疗建议——那属于医生和公共卫生机构的职责，决策应该基于他们的专业评估，而不是一本书。

顺便补上证据史上的第一页：1796 年，英国医生琴纳（E.~Jenner）注意到挤奶女工得过牛痘（一种症状轻微的病）后就不再得致命的天花，于是把牛痘材料接种给一个男孩，之后男孩接触天花也安然无恙。琴纳不知道病毒、不知道抗体，他只掌握了一条经验规律：“得过相近的轻病，能防重病。”但这条经验规律被两百年间亿万人的数据反复验证，最终让人类在 1980 年彻底消灭了天花——第一种被人类主动抹掉的传染病。机制的解释，要等克隆选择学说出现才补上；\textbf{先有有效的经验，后有机制的理论}，这在医学史上是常态，提醒我们“知道有效”和“知道为什么有效”是两份不同的证据。

\section{系统也会出错：过强、打偏、缺失}

任何识别系统都会犯错，免疫系统也不例外。出错的模式恰好有三种，对应“识别—反应—记忆”链条的不同环节，也是本章模型的边界。

\begin{itemize}[itemsep=2pt]
\item \textbf{过敏：对无害目标过度反应}。花粉、尘螨、花生里的某些蛋白质，对大多数人是无害的，但有些人的免疫系统偏偏把它们当成大敌，制造出一类专门引发炎症的抗体（IgE）。下次再遇到，肥大细胞瞬间释放大量组胺，血管通透性暴涨——打喷嚏、流涕、荨麻疹，严重的会发生可致命的全身过敏反应。攻击是真的，只是目标选错了。为什么偏偏是这些人？遗传、幼年环境暴露都有关，“卫生假说”认为过于干净的童年可能让免疫系统“缺乏训练素材”，但这仍是一个\textbf{证据尚在积累}的假说，不是定论。
\item \textbf{自身免疫：把枪口对准自己}。还记得“每个细胞公示内部样品”的检查制度吗？建立这套制度时，本该把那些碰巧能识别自身抗原的淋巴细胞\textbf{删除}——这个审查难免漏网。一旦某个漏网的克隆被激活，免疫系统就开始攻击自己的组织：1 型糖尿病是胰岛细胞被杀伤，类风湿关节炎是关节被持续炎症攻击。目前这类疾病大多只能控制、难以根治，治疗思路往往是抑制免疫——代价是防御力随之下降，两害相权。
\item \textbf{免疫缺陷：防线缺了一块}。天生缺少某类免疫细胞或抗体（原发性免疫缺陷），或者免疫系统被病毒摧毁——HIV 专门攻击辅助 T 细胞，等于炸掉指挥部，整个免疫网络随之瘫痪，这就是艾滋病致死的原因：死于各种平时无害的微生物的轮番感染。
\end{itemize}

把这三类摆在一起，你会看到一个教科书上不常明说的事实：免疫系统的目标从来不是“越强越好”，而是\textbf{精确}——在“漏掉敌人”和“误伤自己”之间走钢丝。进化筛选出的不是一个火力最猛的系统，而是一个犯错代价最小的系统。

\begin{wuqu}
“免疫力越强越好，所以要多吃补品提高免疫力。”——双重错误。第一，“强”如果指反应更猛，那过敏和自身免疫病就是“强”的反面教材：对花粉反应极强的免疫系统，恰恰是出了故障的免疫系统。第二，目前没有任何一种保健品或“补品”拥有高质量证据，能证明它可以把健康人的免疫功能提升到正常水平之上——这类说法大多停留在细胞实验、动物实验或没有对照的“用户体验”层面，证据等级很低。免疫系统正常工作的真正基础毫无商业神秘感：充足睡眠、均衡营养、规律运动、按时接种。如果有产品宣称能“提高免疫力”，请追问一句：证据是随机对照试验，还是广告？
\end{wuqu}

\section{一场没有终点的军备竞赛}

最后一块拼图：敌人不是静止的靶子。

病毒和细菌也在繁殖，繁殖就有变异，变异里碰巧出现“抗原换个形状”的个体，就能躲过已有的抗体识别——你体内的选择压力，对它们来说就是进化的方向盘。流感病毒每年都在换抗原的“外套”，所以流感疫苗要年年更新；麻疹病毒的抗原几十年几乎不变，所以一针麻疹疫苗管几十年。区别不在疫苗技术，在\textbf{敌人变异的快慢}。

宿主这边也在演化：人群里碰巧携带某种受体变体、能抵抗某种病原体的个体，在大瘟疫中更容易活下来留下后代。镰状细胞贫血的基因在疟疾流行区被保留下来，就是这种拉锯留下的痕迹（第 73 章会再遇到它）。病原体与宿主，互相充当对方的选择压力，已经这样共同演化了上亿年——这套“军备竞赛”的逻辑，第十篇会用整整一篇的篇幅展开。

\section{回到那道小伤口}

现在可以兑现开场的承诺了。被纸划破的手指，第一天发红发热：巡逻的先天免疫细胞发现了从破口涌入的细菌，释放的信号分子让局部血管扩张、增援涌入——红肿热痛不是“伤口恶化的表现”，而是\textbf{战场本身}。随后几天，吞噬细胞清理细菌和坏死组织，若有病毒参与，适应性免疫的克隆选择也在淋巴器官里悄悄进行。战斗结束，表皮细胞增殖补上缺口，结痂脱落。如果那道伤口里混进了你曾经接种疫苗针对的病原体，记忆细胞会在你毫无察觉的情况下就把它处理掉——你的接种本之所以管用，靠的不是本子里记录的药液，而是训练后留在你淋巴结和骨髓里的那批记忆细胞。

而针眼大小的疫苗注射处，一场同样微型的“军事演习”完成了同样的事：一次红肿，一批记忆。

\section{下一章的悬念}

不过，本章的模型藏着一个巨大的漏洞，你可能已经察觉了：你的肠道里住着\textbf{上万亿个细菌}，总重量超过一公斤——按“非我就攻击”的规则，它们每一个都该被免疫系统当成敌人消灭。可身体不但容忍它们，离了它们还会出各种问题。免疫系统是怎么学会“对有些非我睁一只眼闭一只眼”的？我们又凭什么和一群微生物共存、甚至互相依赖？答案是：你根本不是一个“个体”，你是一个行走的生态系统。下一章，我们去见见你的房客们。

\begin{tiaozhan}
抗体多样性到底怎么“洗牌”洗出来的？编码抗体可变区的基因，在基因组里存着好几组片段库：V 片段约 $40$ 种、D 片段约 $25$ 种、J 片段约 $6$ 种（以重链为例，数字因统计口径略有出入）。每个 B 细胞成熟时，从每组里随机抽一段拼起来，像从三摞牌里各抽一张：组合数约为 $40\times 25\times 6=6000$ 种重链。轻链再独立地拼一次，贡献约 $300$ 种组合。重链与轻链自由搭配：$6000\times 300\approx 2\times 10^{6}$。但这远不是终点——拼接的“接缝”处还会随机插入或删去几个核苷酸，这一步的随机性没有上限，直接把多样性推到 $10^{11}$ 量级以上。再加上 B 细胞激活后还会在可变区“定点加速突变、择优留下”（体细胞高频突变），最终能识别的形状数以亿亿计。请留意这个设计的深意：多样性来自\textbf{受控的随机}，而不是来自对未来的预言。
\end{tiaozhan}

\section*{练一练}

\begin{sikaoti}
\item 画一张“敌人入侵信息流图”：从皮肤破口开始，画出先天免疫（吞噬细胞）和适应性免疫（B 细胞、T 细胞）各自在什么位置介入，用箭头标出谁给谁传递“警报/确认”信号。图旁注明：箭头表示信息或物质的传递，不代表细胞真的排着队。
\item 用克隆选择的三步逻辑，解释下面两个现象为什么不同：第一次接种某疫苗后一两周才有充足抗体；而半年后再遇到真病原体，抗体一两天内就大量出现。请在你的图上标出“记忆细胞”改变了哪一步。
\item 有人说：“疫苗里的抗原既然能引起免疫反应，说明疫苗本身就是病原体，打了等于生病一次。”请指出这个推理错在哪里（提示：区分“识别目标”与“完整敌人”，并想想减毒、灭活、蛋白片段三类疫苗的差别）。
\item 过敏和自身免疫病常常被笼统地说成“免疫力太强”。根据本章内容，分别说明这两种病真正出问题的环节是什么，并解释为什么“增强免疫力”对它们不是好主意。
\item 查一次资料：找出一种“抗原几乎不变、疫苗一针长效”的病原体和一种“抗原快速变异、疫苗需要更新”的病原体，用“变异为病原体提供多样性、免疫系统提供选择压力”的框架，解释这种差别。这和第 10 章之前学过的哪条演化逻辑是同构的？
\item 估算题：本章算过你体内约有 $2\times 10^{20}$ 个抗体分子。假设一场感染中，记忆 B 细胞扩增出的浆细胞有 $10^{8}$ 个，每个每秒分泌 $2000$ 个抗体，照这个速度，一天能生产多少个抗体分子？和体内的“库存”相比，这个产能说明了什么？
\end{sikaoti}
